\documentclass[a4paper,14pt, unknownkeysallowed]{extreport}
\usepackage{placeins}
\usepackage{cmap} % Улучшенный поиск русских слов в полученном pdf-файле
\usepackage[T2A]{fontenc} % Поддержка русских букв
\usepackage[utf8]{inputenc} % Кодировка utf8
\usepackage[english,russian]{babel} % Языки: русский, английский
%\usepackage{pscyr} % Нормальные шрифты
\usepackage{enumitem}
\usepackage{multirow}
\usepackage[14pt]{extsizes}

\usepackage{caption}
\captionsetup{labelsep=endash}
\captionsetup[figure]{name={Рисунок}}

\usepackage{amsmath}

\usepackage{geometry}
\geometry{left=30mm}
\geometry{right=15mm}
\geometry{top=20mm}
\geometry{bottom=20mm}

\setlength{\parindent}{1.25cm} % для шрифта 14pt

\usepackage{titlesec}
\titleformat{\section}
	{\normalsize\bfseries}
	{\thesection}
	{1em}{}
\titlespacing*{\chapter}{0pt}{-30pt}{8pt}
\titlespacing*{\section}{\parindent}{*4}{*4}
\titlespacing*{\subsection}{\parindent}{*4}{*4}

\usepackage{setspace}
\onehalfspacing % Полуторный интервал

\frenchspacing
\usepackage{indentfirst} % Красная строка

\usepackage{titlesec}
\titleformat{\chapter}{\LARGE\bfseries}{\thechapter}{20pt}{\LARGE\bfseries}
\titleformat{\section}{\Large\bfseries}{\thesection}{20pt}{\Large\bfseries}

\usepackage{listings}
\usepackage{xcolor}

% Для листинга кода:
\lstset{ %
	language=lisp,   					% выбор языка для подсветки	
	basicstyle=\small\sffamily,			% размер и начертание шрифта для подсветки кода
	numbers=left,						% где поставить нумерацию строк (слева\справа)
	%numberstyle=,					% размер шрифта для номеров строк
	stepnumber=1,						% размер шага между двумя номерами строк
	numbersep=5pt,						% как далеко отстоят номера строк от подсвечиваемого кода
	frame=single,						% рисовать рамку вокруг кода
	tabsize=4,							% размер табуляции по умолчанию равен 4 пробелам
	captionpos=t,						% позиция заголовка вверху [t] или внизу [b]
	breaklines=true,					
	breakatwhitespace=true,				% переносить строки только если есть пробел
	escapeinside={\#*}{*)},				% если нужно добавить комментарии в коде
	backgroundcolor=\color{white},
}
\usepackage{pgfplots}
\usetikzlibrary{datavisualization}
\usetikzlibrary{datavisualization.formats.functions}

\usepackage{graphicx}
\newcommand{\img}[3] {
	\begin{figure}[h!]
		\center{\includegraphics[height=#1]{inc/img/#2}}
		\caption{#3}
		\label{img:#2}
	\end{figure}
}
\newcommand{\boximg}[3] {
	\begin{figure}[h]
		\center{\fbox{\includegraphics[height=#1]{inc/img/#2}}}
		\caption{#3}
		\label{img:#2}
	\end{figure}
}

% для матриц
\makeatletter
\renewcommand*\env@matrix[1][\arraystretch]{%
	\edef\arraystretch{#1}%
	\hskip -\arraycolsep
	\let\@ifnextchar\new@ifnextchar
	\array{*\c@MaxMatrixCols c}}
\makeatother

\usepackage{amsfonts}
\usepackage{graphicx}


\usepackage[justification=centering]{caption} % Настройка подписей float объектов

\usepackage[unicode,pdftex]{hyperref} % Ссылки в pdf
\hypersetup{hidelinks}

\usepackage{csvsimple}

\addto\captionsrussian{\def\refname{Список используемых источников}}

\newcommand{\code}[1]{\texttt{#1}}

\makeatletter
\renewcommand{\@biblabel}[1]{#1. \hfill}
\makeatother


\begin{document}


\begin{titlepage}
	\newgeometry{pdftex, left=2cm, right=2cm, top=2.5cm, bottom=2.5cm}
	\fontsize{12pt}{12pt}\selectfont
	\noindent \begin{minipage}{0.15\textwidth}
		\includegraphics[width=\linewidth]{inc/img/b_logo.jpg}
	\end{minipage}
	\noindent\begin{minipage}{0.9\textwidth}\centering
		\textbf{Министерство науки и высшего образования Российской Федерации}\\
		\textbf{Федеральное государственное бюджетное образовательное учреждение высшего образования}\\
		\textbf{«Московский государственный технический университет имени Н.Э.~Баумана}\\
		\textbf{(национальный исследовательский университет)»}\\
		\textbf{(МГТУ им. Н.Э.~Баумана)}
	\end{minipage}
	
	\noindent\rule{18cm}{3pt}
	\newline\newline
	\noindent ФАКУЛЬТЕТ $\underline{\text{«Информатика и системы управления»~~~~~~~~~~~~~~~~~~~~~~~~~~~~~~~~~~~~~~~~~~~~~~~~~~~~~~~}}$ \newline\newline
	\noindent КАФЕДРА $\underline{\text{«Программное обеспечение ЭВМ и информационные технологии»~~~~~~~~~~~~~~~~~~~~~~~}}$\newline\newline\newline\newline\newline\newline\newline
	
	
	\begin{center}
		\Large\textbf{Отчет по лабораторной работе №2}
		\Large\textbf {по дисциплине <<Функциональное и логическое программирование>>}
	\end{center}
	
	\noindent\textbf{Тема} $\underline{\text{~~Определений функций пользователя~~~~~~~~~~~~~~~~~~~~~~~~~~~~~~~~~~~~~~~~~~~~~~~~~~~~~~~~~~~~~~~~~}}$\newline\newline\newline
	\noindent\textbf{Студент} $\underline{\text{~~Котляров Н.А.~~~~~~~~~~~~~~~~~~~~~~~~~~~~~~~~~~~~~~~~~~~~~~~~~~~~~~~~~~~~~~~~~~~~~~~~~~~~~~~~~~~~~~~~~}}$\newline\newline
	\noindent\textbf{Группа} $\underline{\text{~~ИУ7-61Б~~~~~~~~~~~~~~~~~~~~~~~~~~~~~~~~~~~~~~~~~~~~~~~~~~~~~~~~~~~~~~~~~~~~~~~~~~~~~~~~~~~~~~~~~~~~~~~~~}}$\newline\newline
	\noindent\textbf{Оценка (баллы)} $\underline{\text{~~~~~~~~~~~~~~~~~~~~~~~~~~~~~~~~~~~~~~~~~~~~~~~~~~~~~~~~~~~~~~~~~~~~~~~~~~~~~~~~~~~~~~~~~~~~~~~~~~~}}$\newline\newline
	\noindent\textbf{Преподаватель} $\underline{\text{~~Толпинская Н.~Б.~~Строганов~Ю.~В.~~~~~~~~~~~~~~~~~~~~~~~~~~~~~~~~~~~~~~~~~~~~~~~~~~}}$\newline
	
	\begin{center}
		\vfill
		Москва~---~\the\year
		~г.
	\end{center}
	\restoregeometry
\end{titlepage}



\section*{Задание 1}

\textbf{Написать функцию, которая принимает целое число и возвращает первое четное число, не меньшее аргумента.}

\begin{lstlisting}
(defun make_even (x)
	( 	
		if (evenp x)
		x
		(+ x 1)
	)
)
\end{lstlisting}

\section*{Задание 2}

\textbf{Написать функцию, которая принимает число и возвращает число того же знака, но с модулем на 1 больше модуля аргумента.}

\begin{lstlisting}
(defun abs_plus_one (x)
	(	
		if (< x 0)
		(- x 1)
		(+ x 1)
	)
)
\end{lstlisting}

\clearpage

\section*{Задание 3}

\textbf{Написать функцию, которая принимает два числа и возвращает список из этих чисел, расположенный по возрастанию.}

\begin{lstlisting}
(defun make_increase_list (a b)
	(
		if (< a b)
		(cons a (cons b Nil))
		(cons b (cons a Nil))
	)
)
\end{lstlisting}

\section*{Задание 4}

\textbf{Написать функцию, которая принимает три числа и возвращает Т только тогда, когда первое число расположено между вторым и третьим.}

\begin{lstlisting}
(defun is_first_between (first second third)
	(
		if 
			(
				or 
				(and (> first second) (< first third))
				(and (> first third) (< first second))
			)
		t
		nil
	)
)
\end{lstlisting}

\FloatBarrier
\section*{Задание 5}
\textbf{Каков результат вычисления следующих выражений?}

\begin{table}
	\begin{center}
		\begin{tabular}{|l|l|}
			\hline
			(and 'fee 'fie 'foe) & foe \\
			\hline
			(or nil 'fie 'foe) & fie \\
			\hline
			(and (equal 'abc 'abc) 'yes) & yes \\
			\hline
			(or 'fee 'fie 'foe) & fee \\
			\hline
			(and nil 'fie 'foe) & Nil \\
			\hline
			(or (equal 'abc 'abc) 'yes) & T \\
			\hline
		\end{tabular}
	\end{center}
\end{table}

\FloatBarrier

\section*{Задание 6}

\textbf{Написать предикат, который принимает два числа-аргумента и возвращает Т, если первое число не меньше второго.}

\begin{lstlisting}
(defun not_less (x y)
	(>= x y)
)
\end{lstlisting}

\section*{Задание 7}

\textbf{Какой из следующих двух вариантов предиката ошибочен и почему?}

\begin{enumerate}[label=\arabic*)]
	\item (defun pred1 (x) (and (numberp x) (plusp x)))
	\item (defun pred2 (x) (and (plusp x)(numberp x)))
\end{enumerate}

Вариант 2 ошибочен: в случае, если x не является числом, вызов функции plusp вызовет в программе ошибку.

\section*{Задание 8}
\textbf{Решить задачу 4, используя для ее решения конструкции: только IF, только COND, только AND/OR.}

~~

~~

~~

\textbf{IF:}

\begin{lstlisting}
(defun is_first_between (first second third)
	(
		if (> first second) 
		(
			if (< first third)
			t
			nil
		)
		
		(
			if (< first second) 
			(
				if (> first third)
				t
				nil
			)
		)
	)
)
\end{lstlisting}

\textbf{COND:}

\begin{lstlisting}
(defun is_first_between (first second third)
	(
		cond 
		(
			(> first second) 
			(
				cond ((< first third) t)
				     (T nil)
			)
		)
		
		(
			(< first second) 
			(
				cond ((> first third) t)
					 (T nil)
			)
		)
	))
\end{lstlisting}

\textbf{AND/OR:}

\begin{lstlisting}
(defun is_first_between (first second third)
	(
		or 
		(and (> first second) (< first third))
		(and (> first third) (< first second))
	)
)
\end{lstlisting}

\section*{Задание 9}

\textbf{Переписать функцию how-alike, приведенную в лекции и использующую COND, используя только конструкции IF, AND/OR.}

\textbf{IF:}

\begin{lstlisting}
(Defun how_alike (x y)
	(
		if (= x y) 'the_same
		(
			if (equal x y) 'the_same
			(
				if (oddp x)
				(if (oddp y) 'both_odd 'difference)
				(if (evenp y) 'both_even 'difference)
			)
		)
	) 
)
\end{lstlisting}


\textbf{AND/OR:}

\begin{lstlisting}
(Defun how_alike (x y)
	(
		or 
		(and (or (= x y) (equal x y)) 'the_same)
		(and (oddp x) (oddp y) 'both_odd)
		(and (evenp x) (evenp y) 'both_even)
		'difference 
	) 
)
\end{lstlisting}

\end{document}